\documentclass[11pt, a4paper]{article}

\usepackage[T1]{fontenc}
\usepackage[utf8]{inputenc}
\usepackage{tgpagella}
\usepackage{tgcursor}
\usepackage{microtype}
\usepackage[spanish, es-noshorthands]{babel}
\usepackage[table, xcdraw]{xcolor}
\usepackage{booktabs}
\usepackage[most]{tcolorbox}
\usepackage[margin=2.5cm]{geometry}
\usepackage{fancyhdr}

\pagestyle{fancy}
\fancyhf{}
\setlength{\headheight}{16pt}
\lhead{\textbf{UCB Sede Tarija} | Ing. Mecatrónica}
\chead{Guía de Desarrollo: Hito 2}
\rhead{Semana 9 - Sesión 1}

\definecolor{ucbblue}{RGB}{0, 51, 102}

\begin{document}

\begin{center}
    \Large\textbf{\textcolor{ucbblue}{Guía de Requerimientos y Desarrollo: Hito 2}}\\[0.2cm]
\end{center}

\begin{tcolorbox}[colback=yellow!10, colframe=black, title=\textbf{Declaración de Arquitectura de Evaluación}]
    El Hito 2 evalúa la integración de la cinemática analítica, evaluación local y generación de trayectorias, hacia la dinámica. Es el núcleo de la ingeniería de manipuladores[cite: 13].
\end{tcolorbox}

\section*{1. Evidencia Final Requerida por Componente[cite: 13]}
\begin{itemize}
    \item \textbf{IK Analítica (Pieper):} 8 ramas analíticas; lógica de filtrado de ramas; validación obligatoria a través de Cinemática Directa (FK).
    \item \textbf{Jacobiano:} Implementación matemática; al menos un caso de validación manual; interpretación física de la matriz.
    \item \textbf{Singularidades:} Experimento seleccionado documentado; evidencia de rango/determinante donde aplique; interpretación de la movilidad perdida; trazado cerca a singularidad.
    \item \textbf{Generación LSPB:} Ecuaciones de la trayectoria de tiempo $q(t), \dot{q}(t), \ddot{q}(t)$; justificación de parámetros; comprobaciones de continuidad y límites.
    \item \textbf{Dinámica (Pendiente W09-S02):} Modelo dinámico; cálculo de torque $\tau(t)$ a lo largo de la trayectoria generada; interpretación.
    \item \textbf{Notebook:} Ecuaciones, código limpio, manejo explícito de unidades, trazado de gráficas (plots), discusión y conclusiones.
\end{itemize}

\section*{2. ¿Qué puede comenzar AHORA?[cite: 13]}
No espere a la sesión de Dinámica para avanzar. Puede implementar ya mismo:
\begin{itemize}
    \item Esqueleto estructural del Notebook.
    \item Implementación completa del algoritmo de Pieper y verificación FK.
    \item Implementación del Jacobiano.
    \item Esqueleto (Scaffold) del análisis de singularidad.
    \item Funciones de trayectoria LSPB.
    \item Utilidades para graficar (Matplotlib).
\end{itemize}

\section*{3. ¿Qué debe esperar a W09-S02?[cite: 13]}
\begin{itemize}
    \item La matriz de Inercia, Coriolis y Gravedad.
    \item La inyección de $\ddot{q}(t)$ y $\dot{q}(t)$ al modelo para obtener torque.
    \item Discusión integrada final.
\end{itemize}

\end{document}
